\documentclass[11pt,a4paper]{article}
\usepackage[utf8]{inputenc}
\usepackage[T1]{fontenc}
\usepackage{geometry}
\geometry{left=1in, right=1in, top=1in, bottom=1in}

\usepackage{enumitem}
\setlist[itemize]{leftmargin=1.5em, itemsep=2pt, parsep=0pt, topsep=2pt}

\usepackage{titlesec}
\titleformat{\section}{\Large\bfseries\vspace{1ex}}{}{0em}{}[\titlerule]
\titlespacing{\section}{0pt}{2ex}{1ex}

\usepackage{xcolor}
\usepackage{hyperref}
\hypersetup{
    colorlinks=true,
    linkcolor=blue,
    filecolor=magenta,      
    urlcolor=teal,
}

\setlength{\parindent}{0pt}

\begin{document}

\begin{center}
    {\Huge \textbf{Sungjin Nam}} \\ \vspace{4pt}
    Fort Collins, CO $\cdot$ \href{mailto:namsungjinn@gmail.com}{namsungjinn@gmail.com} \\
    \href{https://scholar.google.com/citations?user=aPLJxRwAAAAJ}{Google Scholar} 
    $\cdot$ \href{https://www.linkedin.com/in/sungjin-nam-661a1347/}{LinkedIn}
    % $\cdot$ \href{https://github.com/sungjinnam}{GitHub} 
    \\ \vspace{6pt}
    {\large \textbf{Applied Research Scientist}} \\
    % \textit{Understand and Support Users from Linguistic and Behavioral Data}
\end{center}

\vspace{0.5ex}

Experienced applied research scientist leveraging natural language processing (NLP) and machine learning (ML) techniques to build automated operational solutions and human-centered applications.
Fluently delivers the message with quantitative analysis, visualizations, and prototyping in cross-functional team settings.

\section*{Professional Experience}

\textbf{ACT Education Corp.}, Iowa City, IA \hfill \textit{July 2021 -- Present} \\
\textit{Senior Research Scientist}
\begin{itemize}
    \item Enhancing CRASE, an automated essay scoring system for high-stakes writing assessments, saving $\sim\$1.5$M annually in human scoring costs with ML/LLM modules for main scoring engine, single-score scoring logic, a disturbing-content classifer, and an off-topic detector.
    \item Researching methods for more efficient and effective assessment development and delivery, including LLM-based item parameter prediction models and synthetic data generation for training scoring models.
    \item Designing and prototyping LLM-based tools for content developers and students, including passage-based question generation and essay tutoring applications.
    \item Developing various internal automation solutions, including form assembly, sales region recommendation, open-ended survey response analysis, and question-solving-step generation.
    \item Maintaining and improving existing ML modules in production.
\end{itemize}

\vspace{1ex}

\textbf{PEARSON}, Centennial, CO \hfill \textit{June 2020 -- May 2021} \\
\textit{Data Scientist, Intelligent Systems Research and Development (Aida Calculus product)}
\begin{itemize}
    \item Created neural-network models to represent mathematical equations and find similar items or matching templates.
    \item Developed an IRT-based Bayesian network learner model to estimate students' proficiency.
    \item Implemented RESTful APIs for the app's popular features, such as practice, diagnostic dashboard, and hint selection.
\end{itemize}

\section*{Technical Skills}

\textbf{Programming Skills}
\begin{itemize}
    \item \textbf{Python}: PyTorch, scikit-learn, pandas, NumPy, Matplotlib/seaborn, NetworkX, Flask.
    \item \textbf{R}: LME4, tm, topicmodels, bnlearn, JAGS, ggplot, Shiny.
    \item \textbf{Others}: Databricks, Sagemaker, Docker, AWS, GCP, Bash shell, Git, SQL.
\end{itemize}

\textbf{Research Methods}
\begin{itemize}
    \item Deep neural networks, fine-tuning language models, automated prompt engineering, contrastive learning, synthetic data generation, reinforcement learning, mixed-effect analysis.
    \item Visualization, A/B testing, eye-tracking, mixed research methods.

    % \item \textbf{ML}: Deep neural networks, synthetic data generation, reinforcement learning, active learning, curriculum learning, mixed-effect analysis.
    % \item \textbf{NLP}: Fine-tuning language models, automated prompt engineering, contrastive learning.
    % \item \textbf{Others}: Visualization, A/B testing, eye-tracking, mixed research methods.
\end{itemize}

\section*{Education}

\textbf{Ph.D., Information}, University of Michigan, Ann Arbor, MI \hfill \textit{2020} \\
Thesis: \href{https://deepblue.lib.umich.edu/bitstreams/53077431-bd06-4137-a7e9-f757bb7cfa28/download}{Understanding Language Learning via Machine Learning Methods on Behavioral and Linguistic Data} \\
Advisor: Prof. Kevyn Collins-Thompson

\textbf{M.Sc., Information}, University of Michigan, Ann Arbor, MI \hfill \textit{2014} \\
Specialization in Information Analysis \& Retrieval and Human-Computer Interaction

\textbf{B.A., Psychology and Information Science \& Culture}, Seoul National University \hfill \textit{2012}

\section*{Publications}

\subsection*{Peer-Reviewed Proceedings and Journals}
\begin{itemize}
    \item Demirtaş, M. A., \textbf{Nam, S.}, \& Griffin, G. (2026). Leveraging Human-AI Collaboration for a Passage-Based Question Authoring Tool. In International Conference on Artificial Intelligence in Education (pp. 475-489).
    \item \textbf{Nam, S.} (2026). Prompt Optimization with Verifiable Rewards for Synthetic Essay Generation. In International Conference on Artificial Intelligence in Education (pp. 125-139).
    \item \textbf{Nam, S.} (2025). Domain-Adaptive Automated Essay Scoring with Topic Relevance Learning. In Proceedings of the Second Workshop on Automated Evaluation of Learning and Assessment Content co-located with 26th International Conference on Artificial Intelligence in Education (AIED 2025).
    \item \textbf{Nam, S.} (2025). Ordinal Classification for Transformer-Based Automated Essay Scoring Models. In International Conference on Artificial Intelligence in Education (pp. 252-259).
    \item Yang, K. B., \textbf{Nam, S.}, Huang, Y., \& Wood, S. (2024). Rhetor: Providing LLM-Based Feedback for Students' Argumentative Essays. In European Conference on Technology Enhanced Learning (pp. 201-205).
    \item \textbf{Nam, S.}, Collins-Thompson, K., Jurgens, D., \& Tong, X. (2024). Finding Educationally Supportive Contexts for Vocabulary Learning with Attention-Based Models. In Proceedings of the 2024 Joint International Conference on Computational Linguistics, Language Resources and Evaluation (LREC-COLING 2024) (pp. 7286-7295).
    \item Yang, B., \textbf{Nam, S.}, \& Huang, Y. (2023). ``Why My Essay Received a 4?'': A Natural Language Processing Based Argumentative Essay Structure Analysis. In International Conference on Artificial Intelligence in Education (pp. 279-290). %(Acceptance rate: 21.11\%)
    \item \textbf{Nam, S.}, Bylinskii, Z., Tensmeyer, C., Jain, R., Wigington, C., \& Sun, T. (2020). Why are you reading this? Predicting reading goal and familiarity from people's mobile interaction behaviors. Journal of Vision, 20(11), 951-951.
    \item \textbf{Nam, S.}, Samson, P. (2019). Integrating Students' Behavioral Signals and Academic Profiles in Early Warning System. In International Conference on Artificial Intelligence in Education, 345-357. %(Acceptance rate: 25\%)
    \item \textbf{Nam, S.}, Frishkoff, G. A., \& Collins-Thompson, K. (2018). Predicting Students' Disengaged Behaviors in an Online Meaning-Generation Task. In IEEE Transactions on Learning Technologies 11 (3), 362-375.
    \item \textbf{Nam, S.}, Collins-Thompson, K., \& Frishkoff, G. A. (2017). Predicting Short- and Long-Term Student Learning with a Vocabulary Tutoring System via Semantic Features of Partial Word Knowledge. In Educational Data Mining, 80-87. %(Acceptance rate: 25\%)
    \item Waddington, R. J., \textbf{Nam, S.}, Lonn, S., Teasley, S. D. (2016). Improving Early Warning Systems with Categorized Course Resource Usage. In Journal of Learning Analytics 3 (3), 263-290.
    \item \textbf{Nam, S.} (2016). Predicting Off-task Behaviors in an Adaptive Vocabulary Learning System. In Educational Data Mining, 672-674. (Doctoral Consortium)
    \item Oh, J., \textbf{Nam, S.}, \& Lee, J. (2014). Generating highlights automatically from text-reading behaviors on mobile devices. In CHI'14 Extended Abstracts on Human Factors in Computing Systems (pp. 2317-2322).
    \item \textbf{Nam, S.}, Lonn, S., Brown, T., Davis, C. S., \& Koch, D. (2014). Customized course advising investigating engineering student success with incoming profiles and patterns of concurrent course enrollment. In Proceedings of the Fourth International Conference on LAK. ACM, 16-25. %(Acceptance rate: 27\%)
    \item Waddington, R. J., \& \textbf{Nam, S.} (2014). Practice exams make perfect: incorporating course resource use into an early warning system. In Proceedings of the Fourth International Conference on LAK. ACM, 188-192. %(Acceptance rate: 27\%)
\end{itemize}

\subsection*{Book Chapters}
\begin{itemize}
    \item \textbf{Nam, S.}, Collins-Thompson, K., \& Frishkoff, G. A. (2017). Modeling Off-task Behaviors and Learning Outcomes on a Meaning-Generation Task. In Analytics for Learning White Paper: Measurement in digital environments.
    \item Frishkoff, G. A., Collins-Thompson, K., \textbf{Nam, S.}, Hodges, L., \& Crossley, S. (2016). Dynamic Support of Contextual Vocabulary Acquisition for Reading (DSCoVAR): An intelligent tutor for contextual word learning. In Adaptive Educational Technologies for Literacy Instruction. N.Y., NY: Routledge.
\end{itemize}

\subsection*{Preprints and Technical Reports}
\begin{itemize}
    \item Wood, S., Li, D., \& \textbf{Nam, S.}. (2026). Validity Arguments for Using CRASE\textsuperscript{\textregistered} Confidence Scoring on ACT Writing Test Essays. Research Report. R2611.
    \item Wood, S., \textbf{Nam, S.}, \& Li, D. (2026). CRASE5\textsuperscript{\textregistered} for ACT Writing Technical Report. Research Report. R2544.
    \item \textbf{Nam, S.}, Jurgens, D., \& Collins-Thompson, K. (2022). An Attention-Based Model for Predicting Contextual Informativeness and Curriculum Learning Applications. arXiv preprint arXiv:2204.09885.
    \item \textbf{Nam, S.}, Bylinskii, Z., Tensmeyer, C., Wigington, C., Jain, R., \& Sun, T. (2020). Using Behavioral Interactions from a Mobile Device to Classify the Reader's Prior Familiarity and Goal Conditions. arXiv preprint arXiv:2004.12016.
\end{itemize}

\section*{Community Service}
\begin{itemize}
    \item \textbf{Reviewer Service:}
    \begin{itemize}
        \item NLP: EMNLP, ACL, BEA
        \item HCI: ACM CHI, ACII, UMUAI
        \item Educational Technology: AIED, ACM LAK, ACM L@S, BJET, JATT, EM: I.P., AIME-Con.
    \end{itemize}
    \item Student Organizer of Academic Innovation at Michigan Analytics Talk Series \hfill \textit{Jan 2016 -- May 2019}
    \item Local Arrangement Chair for ACM SIGIR \hfill \textit{2018}
\end{itemize}

\section*{Invited Talks}
\begin{itemize}
    \item Machine Learning Research in Educational Technology (October 2023) at the University of Wisconsin, Computer Science Department.
    \item Learning from Implicit Signals in Learning (August 2018) at Seoul National University, Department of Psychology.
    \item Understanding Online Behavior and Performance Signals in Educational Systems (April 2018) at the University of Oregon, Computer and Information Science Department.
    \item Predicting Short- and Long-Term Student Learning with a Vocabulary Tutoring System via Semantic Features of Partial Word Knowledge (June 2017) at Seoul National University, Department of Communication.
    \item Learning Analytics with Mass Data: Analyzing Higher Education's Academic Performance and Decision Making with Data (June 2014) at Seoul National University, Department of Psychology, and Department of Communication.
\end{itemize}

\section*{Public Presentations}
\begin{itemize}
    \item \textbf{Nam, S.} (2026). Synthetic Data Generation for Automated Essay Scoring Models. Presented at National Council on Measurement in Education Annual Meeting, Los Angeles, CA.
    \item \textbf{Nam, S.} (2025). Improving Zero-Shot Classification of Alerting Content Essays with Automatically Generated Descriptions. Presented at National Council on Measurement in Education Annual Meeting, Denver, CO.
    \item \textbf{Nam, S.} (2025). Hybrid Approach for Understanding Educational Surveys Using Generative Model and Text Encoders. Presented at National Council on Measurement in Education Annual Meeting, Denver, CO.
    \item \textbf{Nam, S.}, \& Samson, P. (2018). Mining Students' In- and Out-of-class Behaviors to Create Earlier Warning System. Presented at Michigan Institute for Data Science Annual Data Science Symposium, Ann Arbor, MI.
    \item \textbf{Nam, S.}, \& Collins-Thompson, K. (2018). Toward Understanding Contextual Information in Text. Presented at University of Michigan Data Science Learning Analytics Challenge Symposium, Ann Arbor, MI.
    \item \textbf{Nam, S.}, Collins-Thompson, K., \& Frishkoff, G. A. (2017). Predicting Short- and Long-Term Student Learning with a Vocabulary Tutoring System via Semantic Features of Definition Responses. Presented at University of Michigan Data Science Research Forum, Ann Arbor, MI.
    \item \textbf{Nam, S.}, \& Collins-Thompson, K. (2017). Understanding Questions with Object and Action Terms: Bi-modal topic modeling of student questions from Piazza. Presented at Michigan Institute for Data Science Annual Symposium, Ann Arbor, MI.
    \item \textbf{Nam, S.}, Collins-Thompson, K., Frishkoff, G. A., Bhide, A., Muth, K. B., \& Perfetti, C. (2016). Modeling Real-time Performance on a Meaning-Generation Task. Presented in American Educational Research Association Annual Meeting, Washington, D.C.; Michigan Institute for Computational Discovery \& Engineering Symposium, Ann Arbor, MI.
    \item Oh, J., \textbf{Nam, S.}, \& Lee, J. (2013). Enhancing Effectiveness of Skim Reading with Highlighting Text. Presented in HCI Korea 2013, Seoul, South Korea.
\end{itemize}

\section*{Research Internship and Assistant Experiences}

\textbf{ADOBE}, San Jose, CA \hfill \textit{June -- Nov 2019} \\
\textit{Machine Learning Intern, Document Intelligence Lab (Mentor: Dr. Tong Sun)}
\begin{itemize}
    \item Designed an experiment to predict mobile device readers' prior knowledge levels and goal conditions from scrolling behaviors. The experiment included a crowdsourcing task with a web application, analyzing eye-tracking signals, user interviews, and developing a mixed-effect classifier.
\end{itemize}

\vspace{1ex}

\textbf{ECHO 360}, Ann Arbor, MI \hfill \textit{May -- Aug 2018} \\
\textit{Graduate Research Intern (Mentor: Dr. Perry Samson)}
\begin{itemize}
    \item Developed early warning prediction model from behavioral logs and academic profiles and suggested customized learning strategies for different student groups.
    \item Supervised undergraduate interns in developing software for transcription, text search, and text clustering.
\end{itemize}

\vspace{1ex}

\textbf{UNIVERSITY OF MICHIGAN}, Ann Arbor, MI \hfill \textit{Sep 2012 -- May 2020} \\
\textit{Graduate Research Assistant}
\begin{itemize}
    \item Developed a deep learning model to identify more instructive sentences for a vocabulary learning system and examined how sentences improved efficiency in training NLP models like word2vec or FastText.
    \item Created an interpretable semantic scale for text by projecting biases in semantic differential with word2vec.
    \item Implemented a prediction model for students' disengaged behaviors and intervention moments.
    \item Expanded a topic model with grammar dependency on online forum texts for richer representation.
    \item Analyzed students' successful learning behaviors, like exploring course-taking patterns from >10 years of enrollment data or behavioral logs from multiple STEM courses' learning management systems.
\end{itemize}

\end{document}
